\documentclass[11pt,a4paper]{article}

\usepackage[T1]{fontenc}
\usepackage[utf8]{inputenc}
\usepackage{lmodern}
\usepackage[margin=2.5cm]{geometry}
\usepackage{microtype}
\usepackage{booktabs}
\usepackage{longtable}
\usepackage{array}
\usepackage{xcolor}
\usepackage[hidelinks]{hyperref}

\renewcommand{\arraystretch}{1.18}

\title{French National Assembly Speech Dataset\\Technical Report}
\author{Vincent Garrigue}
\date{June 2026}

\begin{document}

\maketitle

\begin{abstract}
This report documents the construction of a unified French National Assembly
speech dataset covering debates from 1998 to 2026. The final dataset combines a
historical ParisParl corpus derived from official PDF debate reports with recent
DILA open-data XML records. It also relies on structured deputy, mandate,
party, and presidency reference tables in order to bind speeches to stable
political entities. The target output is an upload-ready speech table with one
row per speech and the minimal schema expected by the downstream Policorp /
Kamil structure.
\end{abstract}

\tableofcontents
\clearpage

\section{Objective}

The objective of the project is to produce a clean, documented, and
reproducible dataset of French National Assembly speeches. The final speech
table must cover the period from 1998 to 2026, keep only speech-level rows, and
expose a compact schema suitable for downstream analysis and database import:

\begin{center}
\begin{tabular}{ll}
\toprule
Column & Description \\
\midrule
\texttt{speech\_id} & Sequential speech identifier after chronological sorting. \\
\texttt{speech\_text} & Cleaned speech text. \\
\texttt{speech\_interjection} & Attached parliamentary reaction or interjection, when available. \\
\texttt{speech\_date} & Speech date in \texttt{DD-MM-YYYY} format. \\
\texttt{speaker\_id} & Deputy identifier or presidency identifier. \\
\texttt{country\_id} & Country reference, fixed to \texttt{3} for France. \\
\texttt{term} & French legislature number at the speech date. \\
\bottomrule
\end{tabular}
\end{center}

The final dataframe is named \texttt{df\_1998\_2026} in the exploratory
notebook and is exported to \path{data/final/speeches_1998_2026.csv}.

\section{Dataset Inventory}

The local repository contains raw sources, intermediate transformed files, and
final analysis-ready tables. Table~\ref{tab:inventory} reports the principal
CSV outputs documented in \path{dataset.md}. Counts are based on local file
line counts, excluding the header row.

\begin{table}[h]
\centering
\small
\begin{tabular}{p{0.50\linewidth}r p{0.26\linewidth}}
\toprule
Local file & Records & Role \\
\midrule
\path{data/final/speeches_1998_2026.csv} & 5,557,923 & Final unified speech table. \\
\path{data/final/parisparl_speeches_1998_2010_processed.csv} & 1,424,597 & Historical ParisParl input for the final merge. \\
\path{data/final/parisparl_speeches_1998_2019_processed.csv} & 2,386,200 & Extended processed ParisParl output. \\
\path{data/speech/2011_2026/centralized_speeches.csv} & 3,084,895 & Centralized DILA XML-derived speech table. \\
\path{data/final/politician.csv} & 3,645 & Deputy reference table. \\
\path{data/final/mandate.csv} & 5,193 & Mandate reference table. \\
\path{data/parties/party.csv} & 41 & Party reference table. \\
\path{data/final/presidency.csv} & 9 & Presidency reference table. \\
\bottomrule
\end{tabular}
\caption{Main local dataset files.}
\label{tab:inventory}
\end{table}

\section{Primary Speech Sources}

\subsection{DILA Open Data, 2011--2026}

The recent speech corpus is built from the DILA open-data archive:
\url{https://echanges.dila.gouv.fr/OPENDATA/Debats/AN/}. DILA is the French
government directorate responsible for legal and administrative information.
For recent French National Assembly debates, this is the preferred source
because it provides the original parliamentary XML structure published by the
administration.

The local DILA pipeline has three reproducible stages:

\begin{enumerate}
  \item \textbf{Download and extraction.} \path{converter/download_dila_debats.py}
  reads the DILA index, detects yearly archives, downloads the requested years,
  and extracts XML files into \path{data/speech/2011_2026/raw}. The extraction
  keeps the year/session directory structure and protects against unsafe archive
  paths.

  \item \textbf{XML conversion.} \path{converter/convert.py} parses each
  \texttt{CRI\_*.xml} debate file and writes structured JSON documents under
  \path{data/speech/2011_2026/converted}. The parser preserves the paired
  \texttt{AAA\_*.xml} source name when present and extracts session metadata,
  paragraph identifiers, speaker blocks, speaker links, speaker roles, and the
  nested debate hierarchy.

  \item \textbf{CSV centralization.} \path{converter/json_to_csv.py} flattens
  converted JSON files into \path{data/speech/2011_2026/centralized_speeches.csv}.
  The result contains one row per paragraph and includes raw text, cleaned text,
  topic hierarchy, speaker metadata, interjections, quotes, and session numbers.
\end{enumerate}

The DILA cleaning stage is required because the XML source is parliamentary
markup rather than an analysis-ready table. Paragraphs may include nested
speaker tags, non-breaking spaces, repeated whitespace, procedural titles,
interruptions, and section headers mixed with spoken content. The converter
therefore stores both raw paragraph text and normalized paragraph text, while
separating speaker metadata from the speech content.

\subsection{ParisParl, 1998--2010 for the Final Merge}

The historical speech corpus is based on the ParisParl corpus:
\url{https://zenodo.org/records/3819374}. ParisParl was produced in the
PolMine project from official PDF reports of plenary sessions in the French
National Assembly. It is used to extend the speech coverage before the DILA XML
period.

The original tabular version was no longer available, but the encoded Corpus
Workbench version was recovered and reversed into yearly CSV files stored under:

\begin{center}
\path{data/speech/1996_2019/parisparl_reden_YYYY.csv}
\end{center}

Although the source description covers 1996 to 2019, the local processed files
currently start in 1998. The first two years are documented as missing in the
data quality review.

Cleaning ParisParl is more complex than cleaning DILA because the source is
PDF-derived rather than native structured XML. The extracted text stream mixes
OCR and detokenization artifacts, page-layout noise, table-of-contents lines,
procedural headings, major topics, subtopics, speeches, and interjections.
Numeric strings can represent article numbers or page references rather than
topic numbers, and presidency or procedure rows can be incorrectly interpreted
as speeches if they are not separated.

The main ParisParl restructuring components are:

\begin{itemize}
  \item \path{mine_parisparl.py} and \path{eda/mine_parisparl.ipynb}, used to
  prototype and apply structure recovery rules.
  \item \path{eda/summary_extraction.py}, which extracts the official
  \texttt{SOMMAIRE} zone and uses it as the preferred source for topic order,
  major topics, and subtopics.
  \item \path{eda/title_cleaning.py}, which fixes French detokenization
  artifacts in parliamentary titles.
  \item \path{eda/interjection_filter.py}, which keeps genuine reactions such
  as applause, laughter, protests, and interruptions while filtering structural
  false positives.
  \item \path{eda/speech_grouping.py}, which groups consecutive rows from the
  same speaker, keeps procedure rows separate, and attaches reaction
  interjections to the relevant speech group.
  \item \path{eda/entity_binding.py}, which links cleaned speeches to party and
  politician references.
\end{itemize}

\section{Reference Entity Tables}

\subsection{Deputies and Mandates}

The deputy and mandate tables provide stable political entities for both speech
corpora. Speeches need deputy identifiers, mandate periods, legislature numbers,
and party or group information so that each row can be linked to a person, a
parliamentary term, and a political family.

The primary structured source is the official Assemblee nationale historical
deputy archive:
\url{https://data.assemblee-nationale.fr/acteurs/historique-des-deputes}.
Local JSON files are stored in:

\begin{itemize}
  \item \path{data/deputes/acteur}
  \item \path{data/deputes/organe}
  \item \path{data/deputes/deport}
\end{itemize}

The conversion is implemented in \path{convert_acteurs_to_csv.py}. It extracts
one politician row per \texttt{PA...} identifier and one mandate row per
\texttt{mp\_id + legislature}. The script also extracts birth information,
profession, institutional status, constituency, mandate dates, election dates,
parliamentary groups, committee memberships, source URLs, and a country
reference.

The Sycomore website is used as a complementary source for legislatures 10 to
17: \url{https://www2.assemblee-nationale.fr/sycomore/recherche}. It fills
missing politicians and mandates, especially for older legislatures where the
local structured JSON data can be incomplete.

The merge rules are conservative:

\begin{itemize}
  \item official local Assemblee nationale JSON rows take priority over
  Sycomore rows;
  \item politicians are deduplicated by \texttt{mp\_id};
  \item mandates are deduplicated by \texttt{mp\_id + term};
  \item existing \texttt{PA...} identifiers are preserved;
  \item generated fallback identifiers are recorded in audit reports;
  \item each mandate must point to an existing politician.
\end{itemize}

\subsection{Political Parties}

\path{data/parties/party.csv} is a manually curated reference table used to
standardize party and parliamentary group labels across deputy, mandate, and
speech datasets. It maps unstable historical labels, acronyms, group names, and
aliases to stable internal party identifiers.

This table is used by \path{convert_acteurs_to_csv.py} when converting mandate
group labels to \texttt{party\_id} values and by \path{eda/entity_binding.py}
when binding speech rows to parties. Labels are normalized before matching:
case, accents, punctuation, parenthetical details, and known historical variants
are handled.

Rows that cannot be mapped safely are not guessed. In the deputy pipeline, they
are emitted as empty or \texttt{P000} values and written to
\path{deputes_boostrap/report/unresolved_party_ids.csv} for manual review.

\subsection{Presidency}

\path{data/final/presidency.csv} is used to bind speeches spoken by the chair
of the Assembly when no deputy identifier is available. During the final merge,
rows identified as presidency speeches are matched to a presidency identifier
according to the speech date.

\section{Final Merge for the 1998--2026 Speech Table}

The final upload-ready speech table combines:

\begin{itemize}
  \item \path{data/final/parisparl_speeches_1998_2010_processed.csv}, for
  historical ParisParl speeches from 1998 to 2010;
  \item \path{data/speech/2011_2026/centralized_speeches.csv}, for DILA
  open-data debates from 2011 to 2026;
  \item \path{data/final/presidency.csv}, for presidency identifier binding.
\end{itemize}

The merge logic follows these steps:

\begin{enumerate}
  \item Reduce DILA rows to the columns needed by the final schema:
  date, cleaned text, deputy identifier, speaker label, and interjection.
  \item Drop rows without a speaker, except presidency rows that can be bound
  later through the presidency table.
  \item Rename DILA columns to the final naming convention.
  \item Recover missing DILA dates from source file names when possible.
  \item Read the processed ParisParl table and restrict it to 1998--2010.
  \item Keep only ParisParl rows with \texttt{row\_type == "speech"}.
  \item Rename ParisParl columns to the final naming convention.
  \item Concatenate the normalized ParisParl and DILA speech tables in
  chronological order while preserving source order within each corpus.
  \item Fill missing presidency speaker identifiers from \path{presidency.csv}
  when the row is spoken by the Assembly presidency.
  \item Set \texttt{country\_id = 3} for France.
  \item Compute \texttt{term} from the speech date using the French legislature
  date ranges from the 11th to the 17th legislature.
  \item Generate \texttt{speech\_id} after sorting, starting at 1.
  \item Remove temporary role, parsed-date, and source-order columns.
\end{enumerate}

\section{Quality Assurance}

The project includes targeted tests under \path{tests/} and notebook-level
quality checks for the final merged table. The final notebook asserts that:

\begin{itemize}
  \item final columns match the expected schema and order;
  \item \texttt{speech\_id} is unique and monotonically increasing;
  \item \texttt{speech\_date} is never missing;
  \item \texttt{speaker\_id} is never empty after presidency binding;
  \item \texttt{country\_id} is always \texttt{3};
  \item \texttt{term} is never missing.
\end{itemize}

Rows that still have no speaker identifier after presidency binding are removed
from the upload dataframe because the target structure requires a bound speaker.

The deputy pipeline also writes audit reports under
\path{deputes_boostrap/report}, including new politicians, new mandates,
fallback identifiers, duplicates, and unresolved party identifiers.

\section{Known Limitations}

\begin{itemize}
  \item The ParisParl source is documented as covering 1996--2019, but the
  local processed files currently start in 1998. The first two years remain a
  known coverage gap.
  \item ParisParl is reconstructed from PDF-derived text. Even after cleaning,
  it is more exposed to OCR artifacts, layout noise, and false structural
  positives than the DILA XML source.
  \item Some fields in the politician target schema are intentionally empty
  because they are not available in the local source JSON files:
  \texttt{education\_level}, \texttt{school}, \texttt{academic\_title},
  \texttt{languages}, \texttt{list}, and \texttt{votes}.
  \item Party labels are historically unstable. The curated party table reduces
  this problem, but unresolved labels are kept for manual review instead of
  being guessed automatically.
  \item The final table keeps the compact Kamil structure. Richer source
  metadata such as detailed topic hierarchy, raw paragraph identifiers, and
  source filenames remains available in intermediate files rather than in the
  final upload table.
\end{itemize}

\section{Reproducibility}

The DILA pipeline can be run with:

\begin{verbatim}
python3 main.py download --years 2011-2026
python3 main.py convert
python3 main.py csv
\end{verbatim}

The full DILA pipeline can also be run with:

\begin{verbatim}
python3 main.py all --years 2011-2026
\end{verbatim}

The deputy and mandate pipeline can be run with:

\begin{verbatim}
python3 convert_acteurs_to_csv.py
\end{verbatim}

Useful checks for the deputy conversion include:

\begin{verbatim}
python3 convert_acteurs_to_csv.py --local-only --dry-run
python3 -m py_compile convert_acteurs_to_csv.py
\end{verbatim}

\section{Recommended Next Steps}

\begin{enumerate}
  \item Finalize the documentation review with the intended dataset consumers.
  \item Verify the 1996--1997 ParisParl coverage gap and decide whether the
  final public scope should be declared as 1998--2026.
  \item Review unresolved party mappings and fallback deputy identifiers before
  database publication.
  \item Keep the compact upload table stable, but preserve intermediate files
  for analyses requiring topic hierarchy, source order, raw text, or source
  filenames.
  \item Run the full test suite before pushing the final tables to the target
  database or VM.
\end{enumerate}

\end{document}
